\section{Companion-source and archive audit}
\label{sec:source-audit}

The observations in this section distinguish a mathematical correction to a displayed formula, the quantitative scope of cited inputs, and the exact locations of historical source passages. A source check does not claim an independent proof of an entire cited paper.

\subsection{C15: the primitive-pairing sign in Milne's report}

In \cite[\S3.6, printed and PDF p.~8]{MilneReport2026}, the primitive space is defined by
\[
P^r(A)=\ker\bigl(L^{g-2r+1}:\mathcal A^r(A)\longrightarrow
\mathcal A^{g-r+1}(A)\bigr),
\]
where $L$ is multiplication by the polarization class $\xi$. The displayed pairing $(a,b)\mapsto ab\xi^{g-2r}$ is called positive definite. Direct visual inspection of the supplied PDF confirms that the sign $(-1)^r$ is absent from that display.

\begin{proposition}[C15: correction of the unsigned display]
With the usual intersection normalization, the positive primitive form is
\[
(a,b)\longmapsto(-1)^r ab\xi^{g-2r}.
\]
The unsigned displayed formula is false already for a primitive divisor on a product of two elliptic curves.
\end{proposition}

\begin{proof}
Let $A=E\times E$, let $\theta_1,\theta_2$ be the pullbacks of the origin divisor, and put $\xi=\theta_1+\theta_2$. Then
\[
\theta_1^2=\theta_2^2=0,\qquad \theta_1\theta_2=1.
\]
The first identities follow because the classes are pulled back from curves; the second follows from the transverse intersection of the coordinate divisors. The nonzero class $a=\theta_1-\theta_2$ is primitive, since $a\xi=0$, whereas $a^2=-2$. Hence the unsigned form cannot be positive definite. The usual Hodge--Riemann sign in degree $k=2r$ and type $(r,r)$ is
\[
(-1)^{k(k-1)/2}i^{r-r}=(-1)^{r(2r-1)}=(-1)^r.
\]
Any factorial normalization is a positive scalar and leaves this sign unchanged. The corrected value in the example is two.
\end{proof}

This corrects the formula rather than refuting the report's standard-conjecture consequence. Codimension two is unchanged; in codimension three R21 and R23 already use the positive norm $-\int\tau^2$. The report separates its rationality and CM Hodge inputs and does not claim in \S2.3 that CM Hodge alone has been shown to imply Hodge for every abelian variety. Its canonical/CM lifting discussion in \S1.9 does not supply a bounded-degree effective lifting construction. These source-scope statements are supported by the direct reading recorded in \texttt{research/followups\_audit.md}. The sign correction is not an entry in the Grothendieck transcription ledger.

\subsection{C16: the quantitative scope of van Hoften--Sempliner}

The supplied \cite[Theorem 2.3.4 and proof, pp.~7--8]{VanHoftenSempliner} clears rational Hodge morphisms by integers $m_i$, obtains a surface cycle pairing nontrivially with a one-dimensional eigenspace, and then uses algebraic projectors and a perfect pairing on algebraic classes to conclude algebraicity of that eigenspace. Its complete argument was read as a qualitative deduction from its cited inputs. The passage does not give a numerical bound for the $m_i$, the selected dual cycle, or its degree.

Theorem 4.0.2 there, p.~12, asserts the existence of a compact open level and four Hodge morphisms with nonzero cup product. The proof in \S4.2.8, pp.~19--22, including Lemmas 4.2.9--4.2.10, uses rational approximation to a real splitting isometry, a finite Schwartz function, and a character decomposition to obtain a nonzero term. Those passages do not state numerical bounds for the rational isometry denominators, level, or eventual cycle degrees. Thus C16 is a checked limitation of what these particular statements provide, not a claim that an effective refinement is impossible. The authors' \S1.3 describes the argument as a variant inspired by the OpenAI paper; it is not presented as a line-by-line verification of that paper. The theorem itself remains a supplied input, without an independent re-proof of the theta and Siegel--Weil machinery.

\subsection{Exact archive locators and their relation to the present bounds}

All scan numbers here are one-based PDF pages, including the cover; the archive folio is one less. The original images were read. We do not infer a universal dictionary for the letters' A/B/C/D labels from a modern naming convention.

In \cite{GrothArchive14}, scan9/folio8 discusses cone/effectivity notation, including $\operatorname{Eff}(G)$. Scans38--39/folios37--38 contain \S3 on $\ell$-adic maps associated with abelian families of cycles and Theorem3.1, with a question mark and a marginal caution about the converse. The annotated passage is not imported as an unconditional biconditional or a degree estimate.

Scans44--47/folios43--46 are the letter to Serre dated27 August1965. Scan44 defines $C^i(X)$ as cycles modulo algebraic equivalence tensored with $\Q$, before formulating its Lefschetz assertion. Scan45 distinguishes a homological reformulation, corrects an earlier claimed equivalence with the K\"unneth statement, and includes the sign $(-1)^m$ in primitive positivity for dimension $2m$. Scans46--47 discuss conditional semisimplicity, correspondence images, affine complements, and stronger equivalence assertions. Our boundedness theorems concern equality of cohomology classes and do not prove those algebraic-equivalence assertions.

In \cite{GrothArchive16}, scans44--47/folios43--46 form the Coates letter dated4 January1967. It discusses the index relation together with Lefschetz-type algebraicity, products with hyperplane sections, specialization, and inverse correspondences. Scans48--50/folios47--49 begin a separate Coates letter, whose handwritten date reads6 January1966. This date is reported as written, without emendation; it corrects our initial grouping of all seven pages under the first letter.

In the second letter, scans48--49 use Cayley--Hamilton to construct an algebraic inverse to a cohomological isomorphism under the stated algebraicity and rational-characteristic-polynomial hypotheses. Scan49 then distinguishes a general monomorphism splitting and expressly invokes the index relation for the additional left-inverse and semisimplicity assertions. Scan50 discusses weak-Lefschetz variants. The qualitative inverse construction does not assert preservation of effective cycles and does not bound a correspondence denominator or degree. Our explicit translated intersections, projectors and charged isogenies supply the separate quantitative statements where claimed.

The corrected letter grouping is a source-locator correction to our own memo. No earlier user transcription with a demonstrable mathematical error was available, so \texttt{TRANSCRIPTION\_NOTES.tsv} remains header-only. The full image-reading record and the separate source corrections are in \texttt{research/archive\_reading.md} and \texttt{sources\_audit/SOURCE\_CORRECTIONS.md}. The actual initial reading and the later refined-locator event are both retained in the check log.
